\chapter{planeteamiento del problema}

El manejo inadecuado de residuos agrícolas representa un desafío crítico para la sostenibilidad ambiental debido a su contribución significativa a la emisión de gases de efecto invernadero (GEI), como el metano (CH\textsubscript{4}) y el dióxido de carbono (CO\textsubscript{2}), producto de la descomposición anaeróbica de desechos orgánicos. Además de las emisiones de GEI, esta gestión deficiente provoca contaminación del suelo y cuerpos de agua, afectando la biodiversidad y generando riesgos para la salud humana \parencite{FAO2021,IPCC2014}. 

Aunque existen estrategias como el compostaje y la disposición controlada de residuos, estas prácticas resultan costosas, ineficientes y de difícil implementación, especialmente en zonas rurales donde los recursos tecnológicos y financieros son limitados \parencite{Smith2014}. Esta situación evidencia la necesidad de soluciones sostenibles, accesibles y eficaces que no solo mitiguen las emisiones de GEI, sino que también optimicen la gestión de residuos agrícolas.

La carencia de sistemas integrados que permitan la medición, reporte, validación (MRV) y trazabilidad de GEI y de variables ambientales claves limita la capacidad del sector agrícola para implementar estrategias efectivas de gestión de residuos. Sin herramientas que proporcionen datos precisos y en tiempo real, resulta complejo tomar decisiones informadas para reducir el impacto ambiental. La implementación de tecnologías emergentes, como el Internet de las Cosas (IoT) y algoritmos de aprendizaje automático (AAA), ofrece la posibilidad de superar estas limitaciones al facilitar el monitoreo continuo y automatizado de variables críticas como temperatura, humedad, pH y concentraciones de GEI \parencite{kamilaris2017review}. Sin embargo, estas tecnologías aún no se han aplicado de forma integral en procesos de gestión de residuos agrícolas, lo que limita su potencial para contribuir a la neutralidad de carbono. 

Una alternativa innovadora es la combinación de IoT y AAA con procesos de bioconversión de residuos que involucran la mosca soldado negra (\textit{Hermetia illucens}), la cual ha demostrado ser altamente eficiente en la transformación de residuos orgánicos en proteína y frass (larvicompost) \parencite{barragan2017nutritional,Bermudez2023}. La integración de sensores IoT permite recopilar datos en tiempo real sobre las condiciones ambientales, mientras que los algoritmos de AAA optimizan las condiciones de biodegradación para maximizar la eficiencia del proceso. Este enfoque no solo mejora la degradación de residuos y reduce las emisiones de GEI, sino que también impulsa la economía circular al generar productos de alto valor comercial, como insumos para alimentos para animales, compuestos bioactivos y bioproductos biodegradables \parencite{Diener2011,Salam2022}.

Frente a este panorama, surge la necesidad de diseñar e implementar un sistema integral que combine IoT y AAA para para la medición, reporte, validación y trazabilidad de GEI y variables ambientales en procesos de gestión de residuos agrícolas. Este sistema permitirá optimizar los procesos de bioconversión, mejorar la precisión de las mediciones y garantizar la trazabilidad de las emisiones, contribuyendo así a la mitigación del cambio climático y a la sostenibilidad del sector agrícola.

Bajo esta premisa, la presente investigación La presente investigación busca desarrollar esta solución tecnológica innovadora, con el propósito de transformar la gestión de residuos agrícolas hacia un modelo más eficiente, sostenible y alineado con los objetivos de neutralidad de carbono.